\section{Fuentes, conservación y reproducción}
El suplemento digital conserva el documento acumulativo íntegro, sus
treinta y nueve fuentes, los programas de comprobación y los recibos
correspondientes. La narrativa aprobada y su continuación se mantienen
literalmente en ese archivo. La presente edición transforma su
enunciación conversacional en exposición académica y registra la
correspondencia entre secciones, sin sustituir los originales.

Los antecedentes proceden de las fuentes de la serie del 21 de
septiembre de 2026. El núcleo formal se incorpora desde su clausura
documental en el artículo constitutivo; los desarrollos de acción,
incidencia, masa y temperatura se conservan con sus demostraciones.
El manifiesto indica cada archivo, intervalo extraído y huella. Estos
datos documentan procedencia y conservación, mientras las pruebas del
cuerpo fijan el alcance matemático de los resultados.

Los programas de composición de constantes, hibridación y
compatibilidad contrastan identidades exactas y ejemplos finitos.
Los diagnósticos numéricos se distinguen de los controles racionales.
Las pruebas de monotonía, convergencia, inversión y paso al límite
permanecen desarrolladas en el texto: los ejemplos ejecutables
complementan esas demostraciones.

La compilación utiliza LuaLaTeX. Las fuentes incluidas y sus figuras
permiten reproducir el PDF sin depender de los directorios originales
del autor. El archivo de instrucciones del paquete especifica el orden
de compilación y el modo de ejecución de las comprobaciones.
